% 原创教学示意；近邻数与簇数是不同超参数。
% 左图查询点 (15,13) 的三个最近邻为 (13,15)、(17,11)、(20,16)。
\begin{tikzpicture}[figure picture,foundation picture,x=1mm,y=1mm,font=\figurefont\fontsize{8.5}{11}\selectfont]
  \foreach \shift in {0,54}{
    \begin{scope}[xshift=\shift mm]
      \fill[FigurePaper] (0,0) rectangle (46,35);
      \foreach \x/\y in {8/8,13/15,17/11,10/19,20/16}{
        \fill[FigureBlue] (\x,\y) circle (0.9);
      }
      \foreach \x/\y in {29/23,34/29,38/23,30/30,39/31}{
        \fill[FigureRed] (\x-0.8,\y-0.8) rectangle (\x+0.8,\y+0.8);
      }
    \end{scope}
  }
  \node[font=\figurefont\small] at (23,41) {$k$近邻：$k=3$};
  \node[font=\figurefont\small] at (77,41) {K-means：$k=2$};
  \draw[FigureStroke,dashed,line width=1.1pt] (15,13) circle (6.7);
  \fill[FigureInk] (15,14.2)--(16.2,13)--(15,11.8)--(13.8,13)--cycle;
  \begin{scope}[xshift=54mm]
    \fill[FigurePaper,opacity=0.82] (0,0) rectangle (46,35);
    \foreach \x/\y in {13.6/13.8,34/27.2}{
      \draw[FigureStroke,line width=1.1pt] (\x-2,\y)--(\x+2,\y) (\x,\y-2)--(\x,\y+2);
    }
  \end{scope}
\end{tikzpicture}
